% SPDX-FileCopyrightText: 2026 IObundle
%
% SPDX-License-Identifier: GPL-3.0-only

The remaining work for the SoCLinux project focuses on full system integration and the development of specialized IP drivers.

\subsection{Interrupt Management (Milestones 2.g, 2.h, 2.i)}
The primary challenge is to provide reliable interrupt support for peripherals under Linux. The planned approach includes:
\begin{itemize}
    \item \textbf{VexRiscv PLIC Resolution:} Resolving existing issues with the VexRiscv's internal PLIC by migrating to external \texttt{iob\_clint} and \texttt{iob\_plic} cores.
    \item \textbf{Interrupt Automation:} Extending the driver generation script to include IRQ handling in the auto-generated kernel modules.
    \item \textbf{Verification:} Using the \texttt{iob\_timer} as a reference core to verify that interrupts are correctly propagated from the hardware to the Linux userspace.
\end{itemize}

\subsection{IP-Specific Linux Drivers (Milestones 3.a, 3.b, 3.c)}
While the \texttt{iob\_linux\_device\_drivers} module provides a solid foundation for CSR access, specialized drivers are needed for performance and compatibility:
\begin{itemize}
    \item \textbf{IOb-Cache:} A driver will be developed to manage the cache as a peripheral accelerator, including control interfaces and performance monitoring.
    \item \textbf{IOb-UART16550:} The core will be validated against the standard Linux \texttt{8250} driver. This ensures that the hardware is fully compatible with existing software ecosystems.
    \item \textbf{IOb-Eth:} Similar to the UART, the Ethernet core will be validated with the \texttt{ethoc} driver, ensuring high-speed network performance and standard interface support.
\end{itemize}
